% 2368-eqno-fraction-number: the equation number is set in text style and may itself hold a fraction stack beside the display (tex.web §§1194, 1199)
% new versus 774-display-eqno-narrow: 774 numbers with a single letter, while this numbers with a fraction so the number line prints a stacked numerator, rule and denominator
\input prelude
\font\tenrm=cmr10 \font\sevenrm=cmr7 \font\fiverm=cmr5
\font\teni=cmmi10 \font\seveni=cmmi7 \font\fivei=cmmi5
\font\tensy=cmsy10 \font\sevensy=cmsy7 \font\fivesy=cmsy5
\font\tenex=cmex10
\skewchar\teni='177 \skewchar\seveni='177 \skewchar\fivei='177
\skewchar\tensy='60 \skewchar\sevensy='60 \skewchar\fivesy='60
\textfont0=\tenrm \scriptfont0=\sevenrm \scriptscriptfont0=\fiverm
\textfont1=\teni \scriptfont1=\seveni \scriptscriptfont1=\fivei
\textfont2=\tensy \scriptfont2=\sevensy \scriptscriptfont2=\fivesy
\textfont3=\tenex \scriptfont3=\tenex \scriptscriptfont3=\tenex
\mathcode`a="7161 \mathcode`b="7162 \mathcode`c="7163
\hsize=300pt \parindent=0pt
\setbox0=\vbox{$$a\eqno{b\over c}$$}
\message{2368 fraceqno ht=\the\ht0}
\setbox1=\vbox{$$a\eqno c$$}
\message{2368 leteqno ht=\the\ht1}
\lsshipbox0
\lsshipbox1
\end
